\chapter{The Many-Electron Problem}
\label{ch:manybody}

\section{The full Hamiltonian}

A piece of gold is a collection of nuclei (charge $Z$, mass $M$, positions
$\RR_I$) and electrons (positions $\rr_i$, spins $\sigma_i$). Ignoring
relativity for now (it returns in force in Chapter~\ref{ch:soc}), the
non-relativistic Hamiltonian in atomic units is
\begin{equation}
\Ha =
\underbrace{-\sum_i \tfrac12 \nabla_i^2}_{\hat T_e}
\underbrace{-\sum_I \tfrac{1}{2M_I}\nabla_I^2}_{\hat T_N}
\underbrace{-\sum_{i,I} \frac{Z_I}{|\rr_i - \RR_I|}}_{\hat V_{eN}}
+\underbrace{\sum_{i<j} \frac{1}{|\rr_i - \rr_j|}}_{\hat V_{ee}}
+\underbrace{\sum_{I<J} \frac{Z_I Z_J}{|\RR_I - \RR_J|}}_{\hat V_{NN}}.
\label{eq:fullH}
\end{equation}
Every property of gold --- its color, conductivity, malleability, melting
point --- is in principle encoded in the eigenstates of
Eq.~\eqref{eq:fullH}. The entire difficulty of condensed-matter theory is
that we cannot solve it.

\section{Born--Oppenheimer: freezing the nuclei}

A gold nucleus is $\sim 3.6\times 10^5$ times heavier than an electron.
Electrons therefore adjust essentially instantaneously to any nuclear
arrangement. Formally, one writes the total wavefunction as
$\Psi(\{\rr\},\{\RR\}) \approx \chi(\{\RR\})\,\Phi_{\RR}(\{\rr\})$, where
$\Phi_{\RR}$ solves the \emph{electronic} problem at clamped nuclei:
\begin{equation}
\Big[\hat T_e + \hat V_{ee} + \hat V_{eN}(\{\RR\})\Big]\,
\Phi_{\RR}(\{\rr\}) = E_{\mathrm{el}}(\{\RR\})\,\Phi_{\RR}(\{\rr\}).
\label{eq:electronicH}
\end{equation}
The electronic energy $E_{\mathrm{el}}(\{\RR\}) + V_{NN}(\{\RR\})$ then acts
as a potential-energy surface on which the nuclei move; its minimum over the
lattice constant $a$ gives the equilibrium structure, and its curvature gives
the bulk modulus. The small parameter of the Born--Oppenheimer expansion is
$\kappa = (m_e/M)^{1/4} \approx 0.04$; the corrections to the clamped-nucleus
picture appear at higher orders in $\kappa$ and are irrelevant for the
ground-state properties we target. From now on
``the problem'' means Eq.~\eqref{eq:electronicH}: $N$ interacting electrons
in the \emph{external potential}
\begin{equation}
\vext(\rr) = -\sum_I \frac{Z_I}{|\rr - \RR_I|}
\end{equation}
created by the fixed nuclei (later: by pseudopotentials).

\section{Why the wavefunction is a dead end}

The electronic wavefunction $\Phi(\rr_1\sigma_1,\dots,\rr_N\sigma_N)$ is a
function of $3N$ continuous coordinates and $N$ spins, antisymmetric under
exchange of any two electrons (Pauli principle). Suppose you discretized
each coordinate axis with a modest 10 points. For a single primitive cell of
gold with 11 valence electrons, storing $\Phi$ would need $10^{33}$ complex
numbers times $2^{11}$ spin components --- more numbers than atoms in a
kilogram of gold. This ``exponential wall'' (Walter Kohn's phrase) is not
a matter of waiting for better GPUs: the
object itself is too large. Any tractable theory must replace the
wavefunction with a compact descriptor.

\section{The electron density}

The natural candidate is the \textbf{electron density}
\begin{equation}
n(\rr) = N \sum_{\sigma_1\cdots\sigma_N}\int
\big|\Phi(\rr\sigma_1, \rr_2\sigma_2, \dots, \rr_N\sigma_N)\big|^2\,
\dd\rr_2 \cdots \dd\rr_N ,
\qquad \int n(\rr)\,\dd\rr = N,
\label{eq:densitydef}
\end{equation}
the probability per unit volume of finding \emph{any} electron at $\rr$,
regardless of what the others are doing. It is a function of just three
variables no matter how many electrons there are: on a $64^3$ real-space
grid it is $2.6\times 10^5$ numbers --- about two megabytes.

Two contributions to the total energy at clamped nuclei are already exact
functionals of $n$ alone:
\begin{equation}
\braket{\Phi|\hat V_{eN}|\Phi} = \int \vext(\rr)\, n(\rr)\, \dd\rr,
\qquad
V_{NN} = \text{constant for fixed nuclei}.
\end{equation}
The kinetic energy and the electron--electron repulsion are not obviously
functionals of the density --- they depend on correlations between electrons,
i.e.\ on the full $\Phi$. The astonishing content of the Hohenberg--Kohn
theorems (next chapter) is that they are functionals of $n$ \emph{anyway}:
the ground-state density secretly determines everything.

\section{Spin densities and magnetization}

Once spin--orbit coupling enters, spin is no longer a spectator. The complete
one-body descriptor is the $2\times 2$ Hermitian \textbf{spin-density
matrix}
\begin{equation}
n_{\alpha\beta}(\rr) = N\!\!\sum_{\sigma_2\cdots\sigma_N}\!\int
\Phi^*(\rr\beta, \rr_2\sigma_2, \dots)\,
\Phi(\rr\alpha, \rr_2\sigma_2, \dots)\, \dd\rr_2\cdots \dd\rr_N ,
\qquad \alpha,\beta\in\{\uparrow,\downarrow\},
\label{eq:spindensmat}
\end{equation}
which can be decomposed on the identity and the Pauli matrices
$\bm\sigma = (\sigma_x,\sigma_y,\sigma_z)$:
\begin{equation}
n_{\alpha\beta}(\rr) = \tfrac12\Big[n(\rr)\,\delta_{\alpha\beta}
+ \mathbf{m}(\rr)\cdot\bm\sigma_{\alpha\beta}\Big],
\qquad
n = \mathrm{tr}\,\hat n, \quad
\mathbf{m}(\rr) = \mathrm{tr}\,[\bm\sigma\, \hat n(\rr)].
\end{equation}
Here $n(\rr)$ is the total density and $\mathbf{m}(\rr)$ the local
\textbf{magnetization density}, a vector field that may point in different
directions at different points (``noncollinear magnetism''). Bulk gold is
non-magnetic --- time-reversal symmetry will force $\mathbf{m}=0$
(Chapter~\ref{ch:soc}) --- but your code should carry the full
$n_{\alpha\beta}$ machinery, because with spinor wavefunctions it costs
little and keeps the theory honest.

\begin{keybox}{Where we stand}
Problem: $N \sim 10^{23}$ interacting electrons; wavefunction exponentially
large. Strategy: reformulate the ground-state problem exactly in terms of the
density $n(\rr)$ (3 variables) or spin-density matrix $n_{\alpha\beta}(\rr)$
(4 real fields: $n$ and $\mathbf m$), then approximate the one piece that has no known closed
form --- exchange and correlation.
\end{keybox}
